% Gabarit pandoc des corrigés publiés (brut et guidé). Contenu sans marque ni nom de site.
% Compilation : lualatex (latexmk -lualatex).
% Le Markdown reste la source unique : ne pas éditer ce fichier .tex à la main.
\documentclass[11pt,a4paper]{article}

% --- Moteur, polices (Latin Modern) et langue ---------------------------------
\usepackage{iftex}
\ifLuaTeX\else\errmessage{Ce document se compile avec lualatex}\fi
\usepackage{amsmath,amssymb,mathtools}
\usepackage{fontspec}
% mathrm=sym : \mathrm prend ses glyphes dans la police mathématique (ℓ, μ, Ω), comme KaTeX.
\usepackage[math-style=TeX,bold-style=TeX,mathrm=sym]{unicode-math}
% Repli pour les glyphes absents de Latin Modern (indices et exposants Unicode,
% flèches, filets des schémas en texte, symboles) : DejaVu, puis Latin Modern Math
% (ℓ, lettres mathématiques Unicode des légendes).
\directlua{
  luaotfload.add_fallback("repliTexte", {"DejaVuSerif:mode=harf;", "DejaVuSans:mode=harf;", "file:latinmodern-math.otf:mode=harf;"})
  luaotfload.add_fallback("repliMono", {"DejaVuSansMono:mode=harf;", "DejaVuSans:mode=harf;", "file:latinmodern-math.otf:mode=harf;"})
}
\setmainfont{Latin Modern Roman}[RawFeature={fallback=repliTexte}]
\setsansfont{Latin Modern Sans}[RawFeature={fallback=repliTexte}]
\setmonofont{Latin Modern Mono}[RawFeature={fallback=repliMono}]
\setmathfont{Latin Modern Math}
\usepackage[french]{babel}
\frenchsetup{SmallCapsFigTabCaptions=false}

% --- Mise en page --------------------------------------------------------------
\usepackage[a4paper,margin=2.2cm,headheight=14pt,footskip=1.2cm]{geometry}
\usepackage{microtype}
\setlength{\emergencystretch}{4em}
\tolerance=1500
\setlength{\parindent}{0pt}
\setlength{\parskip}{0.55\baselineskip plus 2pt minus 1pt}
\usepackage{xcolor}
% Gris nommés : babel-french rend « ! » actif dans le document.
\definecolor{gris2}{gray}{0.98}
\definecolor{gris3}{gray}{0.97}
\definecolor{gris5}{gray}{0.95}
\definecolor{gris35}{gray}{0.65}
\definecolor{gris40}{gray}{0.60}
\definecolor{gris45}{gray}{0.55}
\definecolor{gris50}{gray}{0.50}
\definecolor{gris55}{gray}{0.45}
\definecolor{gris60}{gray}{0.40}
\definecolor{gris65}{gray}{0.35}
\definecolor{gris70}{gray}{0.30}
\definecolor{gris80}{gray}{0.20}
\definecolor{gris85}{gray}{0.15}
\usepackage{graphicx}
\usepackage{float}
\usepackage{caption}
\captionsetup{font=small,labelformat=empty,justification=centering}
\usepackage{booktabs,longtable,array,multirow,calc}
\usepackage{siunitx}
\sisetup{locale=FR}
\usepackage[version=4]{mhchem}
\usepackage{listings}
\usepackage{fancyhdr}
\usepackage{lastpage}
\usepackage[most]{tcolorbox}
\usepackage{enumitem}
\setlist{nosep,leftmargin=*,topsep=2pt}
\usepackage[hidelinks,unicode]{hyperref}
\usepackage{bookmark}
\hypersetup{pdftitle={Banque PT PT — Langue vivante B 2026},pdfauthor={Rédigé par IA},pdfcreator={LaTeX},pdfsubject={Corrigé guidé}}
\usepackage{xurl}

% --- Pandoc : listes, tableaux, figures ----------------------------------------
\providecommand{\passthrough}[1]{#1}
\providecommand{\tightlist}{\setlength{\itemsep}{0pt}\setlength{\parskip}{0pt}}

\makeatletter
\def\maxwidth{\ifdim\Gin@nat@width>\linewidth\linewidth\else\Gin@nat@width\fi}
\def\maxheight{\ifdim\Gin@nat@height>0.55\textheight 0.55\textheight\else\Gin@nat@height\fi}
\makeatother
\setkeys{Gin}{width=\maxwidth,height=\maxheight,keepaspectratio}
\setcounter{secnumdepth}{-\maxdimen}
\providecommand{\pandocbounded}[1]{#1}
\usepackage{longtable,booktabs,array}

% --- Code (listings) -----------------------------------------------------------
\lstset{
  basicstyle=\ttfamily\small,
  columns=fullflexible,
  keepspaces=true,
  breaklines=true,
  breakatwhitespace=false,
  postbreak=\mbox{\textcolor{gris50}{$\hookrightarrow$}\space},
  frame=single,
  rulecolor=\color{gris35},
  framesep=4pt,
  xleftmargin=5pt,xrightmargin=5pt,
  backgroundcolor=\color{gris3},
  keywordstyle=\bfseries,
  commentstyle=\color{gris65},
  stringstyle=\color{gris80},
  showstringspaces=false,
  tabsize=4,
  extendedchars=true,
  aboveskip=6pt,belowskip=6pt,
}

% --- Macros des corrigés (KaTeX ↔ LaTeX) ---------------------------------------
% Les corrigés sont écrits pour KaTeX : on définit ici ce que KaTeX connaît et
% que LaTeX n'a pas, sans jamais retoucher le Markdown.
\providecommand{\R}{\mathbb{R}}
\providecommand{\N}{\mathbb{N}}
\providecommand{\Z}{\mathbb{Z}}
\providecommand{\Q}{\mathbb{Q}}
\providecommand{\C}{\mathbb{C}}
\providecommand{\dd}{\mathrm{d}}
\providecommand{\llbracket}{\lBrack}
\providecommand{\rrbracket}{\rBrack}
\providecommand{\lBrace}{\lbrace}
\providecommand{\rBrace}{\rbrace}
\providecommand{\Box}{\square}
\providecommand{\Diamond}{\lozenge}
\providecommand{\varnothing}{\emptyset}
\providecommand{\argmax}{\operatorname*{arg\,max}}
\providecommand{\argmin}{\operatorname*{arg\,min}}
\renewcommand{\leq}{\leqslant}
\renewcommand{\geq}{\geqslant}

% Équation hors texte : réduite à la largeur de la ligne si elle déborde.
\newsavebox{\boiteequation}
\newcommand{\equationajustee}[1]{%
  \sbox{\boiteequation}{$\displaystyle #1$}%
  \ifdim\wd\boiteequation>\linewidth
    \[\resizebox{\linewidth}{!}{\usebox{\boiteequation}}\]%
  \else
    \[#1\]%
  \fi}

% --- Encadrés pédagogiques (sobres, lisibles en noir et blanc) -----------------
\tcbset{
  encadre/.style={
    enhanced,breakable,sharp corners,
    colback=white,colframe=gris55,
    boxrule=0.4pt,leftrule=2.2pt,
    left=6pt,right=6pt,top=3pt,bottom=3pt,
    before skip=6pt,after skip=6pt,
    fonttitle=\bfseries\small,coltitle=black,
    attach title to upper,after title={\quad},
    parbox=false,
  }
}
\newtcolorbox{indication}[1]{encadre,colframe=gris45,colback=gris2,title={Indication~#1}}
\newtcolorbox{methode}{encadre,colframe=gris70,title={Méthode}}
\newtcolorbox{piege}{encadre,colframe=gris85,boxrule=0.8pt,leftrule=2.2pt,title={Piège}}
\newtcolorbox{jury}{encadre,colframe=gris60,colback=gris5,title={Jury}}
\newtcolorbox{encadrenote}{encadre,colframe=gris40}

% --- Titres de section ---------------------------------------------------------
\usepackage{titlesec}
\titleformat{\section}{\Large\bfseries}{}{0pt}{}
\titleformat{\subsection}{\large\bfseries}{}{0pt}{}
\titleformat{\subsubsection}{\normalsize\bfseries}{}{0pt}{}
\titlespacing*{\section}{0pt}{14pt plus 4pt minus 2pt}{6pt}
\titlespacing*{\subsection}{0pt}{10pt plus 3pt minus 2pt}{4pt}
\titlespacing*{\subsubsection}{0pt}{8pt plus 2pt minus 2pt}{2pt}

% --- En-tête et pied de page ---------------------------------------------------
\pagestyle{fancy}
\fancyhf{}
\fancyhead[L]{\small\textsc{Corrigé guidé} · Banque PT PT}
\fancyhead[R]{\small Langue vivante B 2026}
\fancyfoot[L]{\footnotesize Corrigé rédigé par IA, non officiel}
\fancyfoot[R]{\footnotesize \thepage\,/\,\pageref*{LastPage}}
\renewcommand{\headrulewidth}{0.4pt}
\renewcommand{\footrulewidth}{0pt}
\fancypagestyle{plain}{\fancyhf{}\fancyfoot[R]{\footnotesize \thepage\,/\,\pageref*{LastPage}}\fancyfoot[L]{\footnotesize Corrigé rédigé par IA, non officiel}\renewcommand{\headrulewidth}{0pt}}

\begin{document}
\thispagestyle{plain}
\begin{center}
  {\small\scshape Corrigé guidé}\par\smallskip
  {\LARGE\bfseries Banque PT PT — Langue vivante B 2026\par}\medskip
  {\large Banque PT · PT · Langue vivante B · 2026\par}
\end{center}
\vspace{2pt}\hrule\vspace{8pt}

\emph{Corrigé rédigé par IA, non officiel. Signalez toute erreur.}

\hypertarget{pruxe9sentation-du-sujet}{%
\section{Présentation du sujet}\label{pruxe9sentation-du-sujet}}

\textbf{Format.} L'épreuve dure 3 heures, sans dictionnaire ni appareil électronique. Le candidat compose dans la langue choisie à l'inscription, faute de quoi il reçoit zéro. Le sujet comporte deux exercices :

\begin{longtable}[]{@{}
  >{\raggedright\arraybackslash}p{(\columnwidth - 6\tabcolsep) * \real{0.1178}}
  >{\raggedright\arraybackslash}p{(\columnwidth - 6\tabcolsep) * \real{0.3629}}
  >{\raggedright\arraybackslash}p{(\columnwidth - 6\tabcolsep) * \real{0.2163}}
  >{\raggedright\arraybackslash}p{(\columnwidth - 6\tabcolsep) * \real{0.2730}}@{}}
\toprule\noalign{}
\begin{minipage}[b]{\linewidth}\raggedright
Exercice
\end{minipage} & \begin{minipage}[b]{\linewidth}\raggedright
Nature
\end{minipage} & \begin{minipage}[b]{\linewidth}\raggedright
Longueur
\end{minipage} & \begin{minipage}[b]{\linewidth}\raggedright
Ce qui est noté (rapport d'anglais, p.~4-5)
\end{minipage} \\
\midrule\noalign{}
\endhead
\bottomrule\noalign{}
\endlastfoot
I & contraction croisée : un texte \textbf{français} à résumer dans la langue choisie & 130 mots ± 10 \% (117 à 143) & langue /10, restitution /5, cohésion /5 \\
II & rédaction (\emph{essay}) : réponse argumentée à une question propre à chaque langue & 200 à 220 mots & langue /12, structure /4, contenu /4 \\
\end{longtable}

Le texte à contracter est le même pour toutes les langues. Cinq langues sont proposées : allemand, anglais, arabe, espagnol, italien. Le candidat n'en traite qu'une ; ce corrigé traite les cinq, soit 10 exercices numérotés de 1 à 10. Le barème détaillé n'est publié que par le jury d'anglais (deux exercices notés chacun sur 20) ; rien n'indique qu'il diffère dans les autres langues, mais rien ne le garantit.

\textbf{Le texte.} Alexandre Gefen, \emph{Nouvelle revue d'esthétique}, n° 33, 2024 : un article dense, de registre universitaire, sur l'art produit par l'IA générative. Il progresse en quatre temps :

\begin{longtable}[]{@{}
  >{\raggedright\arraybackslash}p{(\columnwidth - 6\tabcolsep) * \real{0.0490}}
  >{\raggedright\arraybackslash}p{(\columnwidth - 6\tabcolsep) * \real{0.4723}}
  >{\raggedright\arraybackslash}p{(\columnwidth - 6\tabcolsep) * \real{0.2208}}
  >{\raggedright\arraybackslash}p{(\columnwidth - 6\tabcolsep) * \real{0.2279}}@{}}
\toprule\noalign{}
\begin{minipage}[b]{\linewidth}\raggedright
§
\end{minipage} & \begin{minipage}[b]{\linewidth}\raggedright
Idée directrice
\end{minipage} & \begin{minipage}[b]{\linewidth}\raggedright
À garder
\end{minipage} & \begin{minipage}[b]{\linewidth}\raggedright
À laisser
\end{minipage} \\
\midrule\noalign{}
\endhead
\bottomrule\noalign{}
\endlastfoot
1 & Première réaction : l'IA ne serait qu'un moyen bon marché d'obtenir des illustrations, sans esthétique propre ; au mieux un art algorithmique ancien, au pire une machine à transférer des styles, plagier et coller ; ses images n'étonnent que par la recombinaison de formes connues & le jugement réducteur et ses deux degrés & DALL-E, Midjourney, Allan Kaprow, l'exposition américaine de 2023 \\
2 & Cette analyse a un mérite : les artistes ont toujours rêvé de machines à faire de l'art et de la disparition de l'auteur ; on peut aussi voir l'IA comme une production déléguée, à la manière de l'atelier, loin du culte littéraire de l'auteur & le fantasme ancien, la création déléguée & Swift, Jaquet-Droz, Mallarmé \\
3 & Même si l'IA produit des images attendues, savoir qu'elles viennent d'une machine nous plonge dans la « vallée de l'étrange » : on ne peut plus raisonner sur l'intention de l'auteur ni sur le contexte, et l'on répugne à se laisser émouvoir & le trouble, l'intention introuvable, la réticence à l'émotion & Calvino et son « classicisme », le nom de Mori \\
4 & Que montrent ces images ? Ni le monde (comme la photographie) ni un regard humain, mais des représentations de représentations, organisées par la machine selon ses propres règles ; à l'artiste et au spectateur de savoir que faire de cette « vérité » étrangère à l'individu & la nature des représentations, l'appel final & le détail des métadonnées \\
\end{longtable}

En 130 mots, cela donne environ 30 mots pour le premier paragraphe et 30 à 35 pour chacun des trois autres. L'erreur la plus fréquente consiste à s'attarder sur le premier paragraphe, le plus facile, et à sacrifier les deux derniers, les plus abstraits.

\textbf{La question.} Les cinq questions posent le même problème, avec des nuances : \emph{Inwiefern} et \emph{To what extent} (dans quelle mesure), \emph{En qué medida}, إلى أي حد, mais \emph{In che senso} en italien (en quel sens). Toutes parlent des \textbf{formes d'art au pluriel} et des \textbf{pratiques} : il faut sortir de l'image, sujet du texte, pour parler de musique, de cinéma, de photographie, de jeu vidéo, de littérature.

\textbf{Ce qu'il ne faut pas laisser passer.} Respecter les deux fourchettes et écrire le nombre de mots à la fin de chaque exercice : le jury d'anglais rappelle que « Des pénalités s'appliquent si le nombre de mots demandé n'est pas respecté » (p.~4). Dans la contraction, ne rien ajouter et n'escamoter aucun paragraphe. Dans l'essai, répondre à la question de la « révolution », pas à une autre (avantages et inconvénients de l'IA, remplacement des artistes), avec des exemples précis pris hors du texte.

\textbf{Gestion du temps.} Aucune durée n'est imposée par exercice. Répartition possible : 20 minutes pour lire le texte et en dégager le plan, 60 minutes pour la contraction (brouillon, réduction, mise au propre), 80 minutes pour l'essai, 20 minutes pour relire les accords et compter les mots. Les deux exercices pèsent autant dans le barème publié par le jury d'anglais : l'essai ne doit pas manger le temps de la contraction.

\textbf{Ce que dit le rapport 2026.} Le rapport traite les cinq langues, de façon très inégale : trois pages en allemand, quatre en anglais, une en arabe, une en espagnol, une en italien. Les pages citées sont celles du PDF du rapport (il n'est pas paginé). Chiffres publiés : en anglais, 2 599 candidats et une moyenne de 9,45/20, la contraction (9,31) ayant été moins réussie que l'essai (9,61) ; en espagnol, 23 candidats, notes de 3 à 19, moyenne 9,74 ; en italien, cinq candidats. Les points que les jurys donnent pour \textbf{discriminants} :

\begin{itemize}
\tightlist
\item
  dans la contraction, la restitution des deux derniers paragraphes : le rapport de nature des images et notre « difficile rapport » aux créations de l'IA (allemand, italien, anglais) ;
\item
  en anglais, la traduction de « vallée de l'étrange » par \emph{uncanny valley}, valorisée, et des exemples \textbf{anglophones} dans l'essai ;
\item
  partout, un essai qui ne recopie pas les arguments du texte et qui cite des exemples précis de plusieurs arts.
\end{itemize}

\hypertarget{partie-i-allemand}{%
\section{Partie I --- Allemand}\label{partie-i-allemand}}

Le jury d'allemand juge la méthode acquise mais la restitution « moins bien réussie cette année, souvent incomplète » (p.~1). Il insiste sur les bases : déclinaison, genre, verbes, orthographe, ponctuation.

\hypertarget{question-1-allemand-i.-contraction}{%
\subsection{Question 1 --- Allemand, I. Contraction}\label{question-1-allemand-i.-contraction}}

\begin{indication}{1}

Avant d'écrire, résumez chaque paragraphe du texte français en une phrase française. Le texte a quatre temps ; votre contraction doit en avoir quatre.

\end{indication}

\begin{indication}{2}

Le premier paragraphe expose un jugement que l'auteur ne reprend pas entièrement à son compte : marquez cette distance. Le discours indirect allemand dispose pour cela d'un mode.

\end{indication}

\begin{indication}{3}

Préparez le lexique avant de rédiger : \emph{künstlich} (artificiel) et \emph{künstlerisch} (artistique), \emph{der Künstler} (pluriel \emph{die Künstler}), \emph{der Betrachter}, \emph{das Kunstwerk}, \emph{die Darstellung}. Pour la « vallée de l'étrange », \emph{das unheimliche Tal} (ou \emph{Uncanny Valley}).

\end{indication}

\hypertarget{solution}{%
\subsubsection{Solution}\label{solution}}

Le plan suit le texte : jugement réducteur (deux phrases), mérite de cette analyse (deux phrases), trouble de la « vallée de l'étrange » (une phrase), nature des images et appel final (deux phrases). Les noms propres et les exemples disparaissent.

Angesichts der Bilderflut generativer KI könnte man diese als billiges Illustrationswerkzeug abtun. Es gäbe keine künstliche Ästhetik: Bestenfalls setze KI die algorithmische Kunst fort, schlimmstenfalls übertrage sie nur bestehende Stile, plagiiere und erstelle geschickte Collagen, die bloß durch neu kombinierte bekannte Formen überraschen. Diese Sichtweise erinnert jedoch daran, dass Künstler seit Langem von kunstschaffenden Maschinen und vom Verschwinden des Autors träumen. Zudem erscheint KI als delegierte Produktion, ähnlich der kollektiven Werkstattarbeit. Da wir aber wissen, dass eine Maschine diese Werke geschaffen hat, geraten wir ins „unheimliche Tal``: Wir können keine Absicht rekonstruieren und lassen uns ungern von ihnen berühren. Schließlich zeigen diese Bilder weder die Welt noch einen menschlichen Blick, sondern Darstellungen, die die Maschine nach eigenen Regeln ordnet. Künstler und Betrachter müssen entscheiden, was sie mit dieser fremden Wahrheit anfangen.

\textbf{(131 Wörter)}

\begin{methode}

Le \emph{Konjunktiv I} (\emph{setze}, \emph{übertrage}, \emph{plagiiere}, \emph{erstelle}) signale que l'on rapporte une thèse sans la partager : c'est exactement le statut du premier paragraphe. Un connecteur (\emph{jedoch}, \emph{zudem}, \emph{aber}, \emph{schließlich}) ouvre ensuite chaque temps de l'argumentation.

\end{methode}

\begin{piege}

Après \emph{da}, \emph{dass} et un pronom relatif, le verbe conjugué va en fin de proposition (\emph{dass eine Maschine diese Werke geschaffen \textbf{hat}}). Et la virgule sépare toujours principale et subordonnée, ce que le jury rappelle avec l'exemple \emph{Es ist schwierig, ein Kunstwerk zu verstehen, wenn man mit ihm keine Emotionen verbindet} (p.~3).

\end{piege}

\begin{jury}

Rapport 2026 (p.~1) : les candidats ont en général restitué les idées principales, mais certains avaient du mal à intégrer « l'idée du difficile rapport aux créations de l'IA ainsi que la question de la nature de ces représentations ». Ce sont les paragraphes 3 et 4 : ils occupent ici plus de la moitié de la contraction.

\end{jury}

\hypertarget{question-2-allemand-ii.-ruxe9daction}{%
\subsection{Question 2 --- Allemand, II. Rédaction}\label{question-2-allemand-ii.-ruxe9daction}}

\begin{indication}{1}

\emph{Inwiefern} demande un degré, pas un oui ou un non. Distinguez ce que l'IA change (les outils, la production, le marché) et ce qu'elle ne change pas forcément (le sens, l'intention).

\end{indication}

\begin{indication}{2}

La question parle de \emph{die verschiedenen Kunstformen} : prévoyez au moins deux arts autres que l'image (musique, cinéma, jeu vidéo, photographie). Un exemple allemand ou germanophone est un atout.

\end{indication}

\begin{indication}{3}

Introduction courte (un fait d'actualité, la question reformulée sans en changer le sens), deux parties avec une phrase-thèse chacune, conclusion qui répond nettement : \emph{Meiner Meinung nach\ldots{}}

\end{indication}

\hypertarget{solution-1}{%
\subsubsection{Solution}\label{solution-1}}

Plan : une accroche sur l'affaire Eldagsen (photographe allemand, prix refusé en 2023), puis une première partie sur la transformation réelle des pratiques (la \emph{Beethoven X} achevée par un programme et créée à Bonn en 2021, le cinéma et le jeu vidéo, la peur des métiers), une seconde sur les limites de cette révolution (le précédent de la photographie, l'intention humaine), et une conclusion qui tranche.

Im Jahr 2023 gewann der deutsche Fotograf Boris Eldagsen einen renommierten Fotopreis mit einem Bild, das eine KI erzeugt hatte, und lehnte ihn anschließend ab. Dieser Fall zeigt, wie sehr die künstliche Intelligenz die Kunstwelt verunsichert. Revolutioniert sie aber wirklich die Kunst?

Einerseits verändert die KI die künstlerische Praxis grundlegend. In der Musik wurde 2021 in Bonn Beethovens zehnte Sinfonie uraufgeführt, die ein Team aus Musikwissenschaftlern und Informatikern mithilfe eines Programms aus Skizzen des Komponisten vollendet hatte. Im Film und in Videospielen erzeugen Programme Landschaften, Stimmen und Figuren in wenigen Sekunden, wofür früher ganze Teams monatelang arbeiteten. Dadurch kann jeder, auch ohne Ausbildung, Werke schaffen, und viele Illustratoren, Synchronsprecher und Musiker fürchten um ihre Arbeit.

Andererseits ist diese Revolution eher technisch als künstlerisch. Schon die Fotografie galt im 19. Jahrhundert als Ende der Malerei, doch sie hat die Maler befreit und ist selbst zu einer Kunstform geworden. Außerdem bleibt der Mensch unverzichtbar: Er formuliert den Prompt, wählt aus und gibt dem Werk einen Sinn. Ohne menschliche Absicht ist ein Bild nur eine Kombination von Daten.

Meiner Meinung nach revolutioniert die KI also vor allem die Werkzeuge und den Kunstmarkt, aber nicht das Wesen der Kunst. Die Künstler müssen jedoch lernen, mit ihr zu arbeiten, statt sie zu ignorieren.

\textbf{(208 Wörter)}

\begin{methode}

\emph{Einerseits \ldots{} andererseits}, \emph{außerdem}, \emph{dadurch}, \emph{doch}, \emph{also}, \emph{jedoch} : six connecteurs suffisent à rendre visible la structure. Le jury apprécie « L'intégration quasi systématique de connecteurs » (p.~1).

\end{methode}

\begin{piege}

Le jury relève l'accusatif oublié après les verbes transitifs, la déclinaison après les prépositions mixtes et le \emph{zu} fautif après un verbe de modalité (\emph{Die künstliche Intelligenz darf die Künstler nicht ersetzen}, p.~2). Vérifiez aussi le pluriel invariable des noms en \emph{-er} : \emph{die Künstler}, \emph{die Maler}, \emph{die Illustratoren} (ce dernier en \emph{-en}).

\end{piege}

\begin{jury}

Rapport 2026 (p.~1) : pour ce sujet, il fallait « maîtriser le vocabulaire de base du domaine de la technologie et de l'art ». Le jury cite \emph{der Betrachter}, \emph{die Realität darstellen}, \emph{Daten verarbeiten}, \emph{das Programm}, \emph{der Bildschirm}, et rappelle la différence entre \emph{künstlich} et \emph{künstlerisch}.

\end{jury}

\hypertarget{partie-ii-anglais}{%
\section{Partie II --- Anglais}\label{partie-ii-anglais}}

Le jury d'anglais publie le barème, des statistiques et l'analyse la plus détaillée. La contraction a été plus difficile que d'habitude ; l'essai a été jugé sur la capacité à sortir du texte.

\hypertarget{question-3-anglais-i.-contraction}{%
\subsection{Question 3 --- Anglais, I. Contraction}\label{question-3-anglais-i.-contraction}}

\begin{indication}{1}

Repérez la structure concessive du texte : un jugement réducteur, ses mérites, puis deux objections plus profondes (notre trouble, la nature des images). Vos connecteurs doivent la faire apparaître.

\end{indication}

\begin{indication}{2}

Préparez les équivalents : \emph{vallée de l'étrange} (\emph{uncanny valley}), \emph{plagier} (\emph{plagiarise}), \emph{production déléguée} (\emph{delegated production}), \emph{disparition de l'auteur} (\emph{the author's disappearance}). Gare aux calques signalés par le jury : \emph{disparition}, \emph{automatisation}, \emph{mecanisation}.

\end{indication}

\begin{indication}{3}

Le dernier paragraphe est le plus difficile, et c'est celui que les candidats oublient : il faut au moins dire que ces images ne représentent ni le monde ni un regard humain, et que c'est à l'artiste et au spectateur d'en décider.

\end{indication}

\hypertarget{solution-2}{%
\subsubsection{Solution}\label{solution-2}}

Les quatre temps sont restitués dans l'ordre, avec un connecteur à chaque bascule (\emph{nevertheless}, \emph{yet}, \emph{finally}). Le conditionnel (\emph{would}) signale que la première analyse est une hypothèse que l'auteur dépasse.

Faced with the flood of AI-generated images, one might dismiss generative AI as a cheap illustration tool. There would be no artificial aesthetic: at best, AI would extend algorithmic art; at worst, it would merely transfer existing styles, plagiarise and assemble skilful collages that impress only by recombining familiar forms. This analysis nevertheless recalls that artists long dreamt of art-making machines and of the author's disappearance, and it invites us to see AI as delegated production, like workshop practice. Yet knowing that a machine made these works draws us into the ``uncanny valley'': we cannot reconstruct any intention and are reluctant to be moved. Finally, these images depict neither the world nor a human gaze, but representations reorganised by the machine's own rules. Artists and viewers must decide what to make of this alien truth.

\textbf{(135 words)}

\begin{methode}

Pour réduire sans trahir, on remplace une énumération d'exemples par la catégorie qu'ils illustrent : Swift et Jaquet-Droz deviennent \emph{artists long dreamt of art-making machines} ; DALL-E et Midjourney, \emph{generative AI}.

\end{methode}

\begin{piege}

\emph{AI} ne prend pas d'article et se reprend par \emph{it}, jamais par \emph{he} ou \emph{she} (rapport, p.~6-7). Et après un modal, l'infinitif sans \emph{-s} : \emph{it would transfer}, et non \emph{it would transfers}.

\end{piege}

\begin{jury}

Rapport 2026 (p.~4) : le texte comportait beaucoup d'arguments et « il fallait ici nécessairement faire des choix ». Le jury accepte toute configuration « à partir du moment où aucun passage n'est complètement escamoté », en particulier le dernier paragraphe. La contraction doit « donner un condensé des idées principales » (p.~5), sans multiplier les exemples ni inventer d'idées. Le jury a valorisé les copies qui connaissaient \emph{uncanny valley} (p.~5).

\end{jury}

\hypertarget{question-4-anglais-ii.-essay}{%
\subsection{Question 4 --- Anglais, II. Essay}\label{question-4-anglais-ii.-essay}}

\begin{indication}{1}

Reformulez la question pour vous-même sans la déformer : il s'agit du caractère \textbf{révolutionnaire} de l'IA dans \textbf{plusieurs} arts, pas de ses avantages et inconvénients, ni de la question de savoir si elle remplacera les artistes.

\end{indication}

\begin{indication}{2}

Choisissez deux ou trois exemples anglophones précis et datés, chacun dans un art différent : la musique, les arts visuels, le cinéma.

\end{indication}

\begin{indication}{3}

Deux parties suffisent. Commencez chacune par une \emph{topic sentence} qui annonce l'idée : \emph{Admittedly, AI has transformed\ldots{}} puis \emph{However, this upheaval affects tools and markets more than art itself.}

\end{indication}

\hypertarget{solution-3}{%
\subsubsection{Solution}\label{solution-3}}

Plan : introduction sur la grève de Hollywood de 2023 ; I. l'IA transforme réellement la manière de créer (« Now and Then » des Beatles, récompensé aux Grammy Awards en 2025 ; l'installation \emph{Unsupervised} de Refik Anadol au MoMA) ; II. la révolution touche les outils et les marchés plus que l'art (le précédent de la photographie, le prix refusé par Boris Eldagsen, les métiers menacés) ; conclusion qui répond à \emph{to what extent}.

In 2023, Hollywood screenwriters and actors went on strike partly to protect their work from artificial intelligence. Such fears suggest that AI is overturning artistic practices, but to what extent is it truly revolutionising them?

Admittedly, AI has transformed the way art is made in many fields. In music, machine-learning software isolated John Lennon's voice from an old demo tape, which allowed the Beatles to release ``Now and Then'', awarded a Grammy in 2025. In the visual arts, Refik Anadol's installation Unsupervised, shown at MoMA in New York, turned the museum's own collection into images generated in real time. Never before could a single artist produce so much, so quickly.

However, this upheaval affects tools and markets more than art itself. Photography was once accused of killing painting, yet it eventually freed painters and became an art form in its own right. Similarly, the German photographer Boris Eldagsen, who won a Sony World Photography Award in 2023 with an AI-generated picture, turned down the prize to spark a debate, not to replace photographers. What AI really threatens is the income of illustrators, voice actors and session musicians.

In conclusion, AI is revolutionising artistic practices and the economics of creation, but the meaning of a work still depends on the human intention behind it.

\textbf{(213 words)}

\begin{methode}

La concession structure tout l'essai : \emph{Admittedly \ldots{} However \ldots{} Similarly \ldots{} In conclusion}. Le jury distingue les meilleures copies par l'étendue des structures : passifs, inversions, modaux, concession (rapport, p.~7). L'inversion \emph{Never before could a single artist produce so much} en est un exemple sobre.

\end{methode}

\begin{piege}

Erreurs relevées par le jury : \emph{to what extend} pour \emph{to what extent}, \emph{generativ}, \emph{artistical}, \emph{softwares} et \emph{musics} (indénombrables), \emph{to realize art}, \emph{automatisation}, \emph{inspirated}, \emph{theorical} (rapport, p.~5-6). Le jury met aussi en garde contre le registre familier (\emph{hype}, \emph{vibe}, \emph{level up}).

\end{piege}

\begin{jury}

Rapport 2026 (p.~5) : « Il ne s'agissait pas de parler des avantages ou des inconvénients, mais de discuter de l'aspect révolutionnaire de l'IA sur l'art dans divers domaines. » Les copies qui ont repris l'essentiel des arguments du texte ont été pénalisées. Le jury regrette des exemples majoritairement francophones (la publicité d'Intermarché, très citée) et recommande « d'essayer d'intégrer des exemples issus du monde anglophone ». Le développement tient « plutôt » en deux parties, faute de place pour trois.

\end{jury}

\hypertarget{partie-iii-arabe}{%
\section{Partie III --- Arabe}\label{partie-iii-arabe}}

Le rapport d'arabe tient en une page : niveau contrasté, méthode de la contraction satisfaisante, et deux rappels précis sur le décompte des mots et la paraphrase.

\hypertarget{question-5-arabe-i.-contraction}{%
\subsection{Question 5 --- Arabe, I. Contraction}\label{question-5-arabe-i.-contraction}}

\begin{indication}{1}

Contracter n'est pas paraphraser : on ne suit pas le texte phrase par phrase, on en restitue l'enchaînement d'idées en quatre temps.

\end{indication}

\begin{indication}{2}

Préparez les termes : الذكاء الاصطناعي التوليدي، الفن الخوارزمي، الانتحال، إنتاج مفوض، اختفاء المؤلف. Pour la « vallée de l'étrange », la traduction usuelle est «وادي الغرابة».

\end{indication}

\begin{indication}{3}

Comptez à mesure : chaque mot graphique compte, mais la conjonction و soudée au mot suivant ne compte pas à part.

\end{indication}

\hypertarget{solution-4}{%
\subsubsection{Solution}\label{solution-4}}

Les quatre temps sont marqués par des articulations nettes : فلا \ldots{} غير أن \ldots{} لكن \ldots{} وأخيرا. Aucun nom propre n'est conservé.

أمام الكم الهائل من الصور التي ينتجها الذكاء الاصطناعي التوليدي، قد نختزله في أداة رخيصة لإنتاج الرسوم التوضيحية. فلا وجود في هذه الحال لجمالية اصطناعية: إنه في أحسن الأحوال امتداد للفن الخوارزمي، وفي أسوئها مجرد تقنية تنقل أساليب قائمة وتنتحلها وتركب منها ملصقات بارعة لا تدهشنا إلا بإعادة تركيب أشكال مألوفة. غير أن هذا التحليل يذكرنا بأن الفنانين حلموا طويلا بآلات تنتج الفن وباختفاء المؤلف، ويدعونا إلى النظر إلى الذكاء الاصطناعي بوصفه إنتاجا مفوضا يشبه العمل الجماعي في الورشة. لكن علمنا بأن آلة صنعت هذه الأعمال يوقعنا في «وادي الغرابة»: إذ يتعذر علينا أن نستعيد نية المؤلف أو سياق العمل، ونتحرج من أن نتأثر بها. وأخيرا، لا تصور هذه الصور العالم ولا نظرة إنسانية، بل تمثيلات أعادت الآلة تنظيمها وفق قواعدها الخاصة. ويبقى على الفنان والمتلقي أن يقررا ما يصنعان بهذه الحقيقة الغريبة عنهما.

\textbf{(134 كلمة)}

\begin{methode}

La phrase verbale (verbe en tête) allège la contraction : لا تصور هذه الصور العالم. Le verbe placé avant un sujet pluriel reste au singulier, mais s'accorde en genre.

\end{methode}

\begin{piege}

La dernière phrase de la contraction porte le duel (يقررا، يصنعان، عنهما) : le sujet est double (الفنان والمتلقي). Le passer au pluriel est une faute d'accord.

\end{piege}

\begin{jury}

Rapport 2026 (p.~8) : « Il faut éviter la paraphrase et veiller à ne pas conserver les exemples dans la contraction ». Le jury rappelle aussi que « la conjonction « و » n'est pas un mot et ne doit pas être comptabilisée », et recommande de compter « de manière scrupuleuse le nombre de mots et de l'inscrire à la fin de l'exercice ».

\end{jury}

\hypertarget{question-6-arabe-ii.-ruxe9daction}{%
\subsection{Question 6 --- Arabe, II. Rédaction}\label{question-6-arabe-ii.-ruxe9daction}}

\begin{indication}{1}

Posez une problématique explicite dès l'introduction et annoncez votre réponse : le jury en fait la condition de la cohérence.

\end{indication}

\begin{indication}{2}

Un exemple du monde arabe donne du relief : l'exposition du robot peintre Ai-Da aux pyramides de Gizeh en 2021, les débats sur les voix de chanteurs disparus recréées par l'IA.

\end{indication}

\begin{indication}{3}

Variez les connecteurs (فمن جهة، ومن جهة أخرى، كما، لذلك، وخلاصة القول) : le jury demande d'en soigner le choix, car ils rendent la lecture plus fluide.

\end{indication}

\hypertarget{solution-5}{%
\subsubsection{Solution}\label{solution-5}}

Plan : introduction sur Ai-Da et problématique (révolution des outils, non de l'essence) ; I. la transformation des pratiques (musique, image, cinéma et jeu vidéo, métiers menacés) ; II. une révolution plus technique qu'artistique (le précédent de la photographie, l'intention humaine) ; conclusion.

في عام 2021، عرض الروبوت الرسام «آي-دا» أعماله عند أهرامات الجيزة، بعد أن احتجزته الجمارك المصرية أياما. ويختصر هذا المشهد الطريف سؤالا جديا: إلى أي حد يحدث الذكاء الاصطناعي ثورة في الفن؟ أرى أنه يحدث ثورة حقيقية في أدوات الإبداع وممارساته، لكنه لا يغير جوهر الفن.

فمن جهة، غير الذكاء الاصطناعي طريقة إنتاج الأعمال الفنية في مجالات كثيرة. ففي الموسيقى، صار بإمكان أي شخص أن يؤلف أغنية كاملة في ثوان، بل أن يقلد أصوات مطربين راحلين، وهو ما أثار جدلا واسعا في العالم العربي. كما أصبحت برامج توليد الصور تنتج في لحظات لوحات بأسلوب أي رسام، مما يطرح بحدة مسألة حقوق المؤلف. وفي السينما وألعاب الفيديو، تولد البرامج المناظر والأصوات والشخصيات بسرعة غير مسبوقة. ولذلك يخشى كثير من الرسامين ومؤدي الدبلجة والموسيقيين على مستقبلهم المهني.

ومن جهة أخرى، تبقى هذه الثورة تقنية أكثر منها فنية. فحين ظهر التصوير الفوتوغرافي في القرن التاسع عشر، ظن كثيرون أنه سيقضي على الرسم، لكنه حرره من وظيفة نقل الواقع وفتح الطريق أمام الفن الحديث. كما يظل الإنسان ضروريا، فهو الذي يختار ويصحح ويمنح العمل معناه. فالصورة التي لا تحمل قصدا إنسانيا ليست سوى تركيب من البيانات.

وخلاصة القول، إن الذكاء الاصطناعي يحدث ثورة في الوسائل وفي سوق الفن، أما الفن نفسه فيبقى تعبيرا عن تجربة إنسانية لا تستطيع الآلة أن تعيشها.

\textbf{(207 كلمات)}

\begin{methode}

Pour exprimer le degré demandé par إلى أي حد, opposez deux plans avec أكثر منها : ثورة تقنية أكثر منها فنية. La formule donne la réponse en une ligne.

\end{methode}

\begin{piege}

Le nom qui suit un nombre s'accorde selon le nombre : on écrit 207 كلمات (unités de 3 à 10) mais 134 كلمة (de 11 à 99). L'indication finale du nombre de mots doit elle-même être correcte.

\end{piege}

\begin{jury}

Rapport 2026 (p.~8) : « il est essentiel de poser une problématique explicite dès le départ afin de donner au devoir toute sa cohérence ». Les meilleures copies ont fait la différence par des arguments solides, un fil conducteur logique et un regard critique nuancé.

\end{jury}

\hypertarget{partie-iv-espagnol}{%
\section{Partie IV --- Espagnol}\label{partie-iv-espagnol}}

Le jury d'espagnol juge la méthode acquise et le niveau de langue en progrès, mais reste intransigeant sur la conjugaison et l'accentuation, et signale des contresens sur le texte.

\hypertarget{question-7-espagnol-i.-contraction}{%
\subsection{Question 7 --- Espagnol, I. Contraction}\label{question-7-espagnol-i.-contraction}}

\begin{indication}{1}

Répartissez vos mots \textbf{avant} de rédiger : environ 30 pour le premier paragraphe, 35 pour chacun des autres. C'est précisément le défaut que le jury relève cette année.

\end{indication}

\begin{indication}{2}

Relisez le texte français pour éviter le faux sens : l'auteur n'adhère pas au jugement du premier paragraphe, il le présente comme une « première réaction ».

\end{indication}

\begin{indication}{3}

Préparez \emph{el valle inquietante} (vallée de l'étrange), \emph{la desaparición del autor}, \emph{una producción delegada}, \emph{el trabajo de taller}.

\end{indication}

\hypertarget{solution-6}{%
\subsubsection{Solution}\label{solution-6}}

Le conditionnel (\emph{cabría}, \emph{No habría}) marque l'hypothèse initiale ; \emph{Sin embargo}, \emph{Pero}, \emph{Por último} ouvrent les trois temps suivants, qui occupent les deux tiers de la contraction.

Ante la avalancha de imágenes generadas por IA, cabría ver la IA como una herramienta de ilustración barata. No habría estética artificial: en el mejor caso, arte algorítmico; en el peor, una técnica que transfiere estilos, plagia y fabrica hábiles collages, sorprendentes solo por recombinar formas conocidas. Sin embargo, este análisis recuerda que los artistas soñaron con máquinas creadoras y con la desaparición del autor, e invita a considerarla una producción delegada, semejante al trabajo de taller. Pero saber que una máquina creó estas obras nos sumerge en el «valle inquietante»: no podemos reconstruir ninguna intención y nos resistimos a emocionarnos. Por último, estas imágenes no muestran ni el mundo ni una mirada humana, sino representaciones reorganizadas según las reglas de la máquina. Al artista y al espectador les toca decidir qué hacer con esta verdad ajena.

\textbf{(137 palabras)}

\begin{methode}

\emph{ni \ldots{} ni} (\emph{no muestran ni el mundo ni una mirada humana}) puis \emph{sino} : la double négation suivie d'une rectification restitue en une phrase le dernier paragraphe.

\end{methode}

\begin{piege}

Une subordonnée seule entre deux points n'est pas une phrase : \emph{Porque la IA solo recombina formas conocidas.} est refusé. Et les seules consonnes doublées en espagnol sont celles de CaRoLiNa : \emph{colaborar}, \emph{comunidad}, \emph{imágenes}, jamais \emph{communidad} ni \emph{immágenes}.

\end{piege}

\begin{jury}

Rapport 2026 (p.~9) : certaines copies ont produit « une contraction un peu trop détaillée de la première partie du texte et une contraction bâclée de la deuxième ». Le texte a aussi causé des faux sens, même chez des hispanophones : le jury conseille de lire très régulièrement pour développer leur « compréhension du sens et des nuances ».

\end{jury}

\hypertarget{question-8-espagnol-ii.-ruxe9daction}{%
\subsection{Question 8 --- Espagnol, II. Rédaction}\label{question-8-espagnol-ii.-ruxe9daction}}

\begin{indication}{1}

\emph{En qué medida} : la réponse est un degré. Séparez ce que l'IA transforme et ce qu'elle laisse intact.

\end{indication}

\begin{indication}{2}

Cherchez des exemples du monde hispanique : une chanson qui imitait par IA la voix de Bad Bunny en 2023, le Dalí recréé numériquement par le musée Dalí de Floride en 2019.

\end{indication}

\begin{indication}{3}

Vérifiez chaque forme verbale et chaque accent écrit avant de recopier : c'est la première exigence du jury.

\end{indication}

\hypertarget{solution-7}{%
\subsubsection{Solution}\label{solution-7}}

Plan : introduction sur « NostalgIA » et problématique ; I. la transformation des pratiques (musique, musées, cinéma et jeu vidéo, métiers menacés) ; II. une révolution plus technique qu'artistique (photographie, intention humaine) ; conclusion à la première personne.

En 2023, «NostalgIA», una canción que imitaba con inteligencia artificial la voz de Bad Bunny, se hizo viral y provocó la indignación del cantante puertorriqueño. Este caso plantea una pregunta: ¿está la IA revolucionando realmente el arte o solo sus herramientas?

Por un lado, es innegable que la IA transforma las prácticas artísticas. En la música, cualquiera puede ahora componer una canción completa en pocos segundos. En los museos, el Museo Dalí de Florida «resucitó» al pintor en 2019: un Dalí recreado digitalmente recibe a los visitantes y conversa con ellos. En el cine y en los videojuegos, los decorados y los efectos especiales se generan cada vez más rápido. Por eso, muchos ilustradores, actores de doblaje y músicos temen perder su trabajo.

Por otro lado, esta revolución es más técnica que artística. En el siglo XIX, la fotografía parecía condenar la pintura, pero liberó a los pintores de la obligación de imitar la realidad y abrió el camino al arte moderno. Además, el ser humano sigue siendo indispensable: es él quien elige, corrige y da sentido a la obra. Sin intención, una imagen no es más que una combinación de datos.

En definitiva, creo que la IA está revolucionando los medios y el mercado del arte, pero no su esencia, que sigue siendo humana.

\textbf{(215 palabras)}

\begin{methode}

La préposition \emph{a} devant un complément d'objet humain déterminé : \emph{resucitó al pintor}, \emph{recibe a los visitantes}, \emph{liberó a los pintores}. Le jury la cite parmi les bases à maîtriser.

\end{methode}

\begin{piege}

\emph{Ser} et \emph{estar} : \emph{esta revolución \textbf{es} más técnica que artística} (caractéristique), mais \emph{la IA \textbf{está} revolucionando} (action en cours, périphrase \emph{estar} + gérondif reprise de la question).

\end{piege}

\begin{jury}

Rapport 2026 (p.~9) : « nous n'admettons pas les barbarismes verbaux ! » Pour l'essai, plusieurs candidats « se sont contentés de reprendre ceux du texte », et l'absence ou la trivialité des exemples a été relevée. Enfin, « l'accentuation n'est pas facultative ».

\end{jury}

\hypertarget{partie-v-italien}{%
\section{Partie V --- Italien}\label{partie-v-italien}}

Cinq candidats seulement ont composé en italien ; le jury juge le niveau meilleur que les années précédentes et la contraction bien maîtrisée.

\hypertarget{question-9-italien-i.-contraction}{%
\subsection{Question 9 --- Italien, I. Contraction}\label{question-9-italien-i.-contraction}}

\begin{indication}{1}

Reformulez au lieu de traduire : une traduction littérale du français conduit, selon le jury, à des restitutions grossières ou inexactes.

\end{indication}

\begin{indication}{2}

Repérez l'idée que les candidats ont souvent omise : notre réticence à être ému par une œuvre produite par une machine (paragraphe 3).

\end{indication}

\begin{indication}{3}

Préparez \emph{la valle perturbante} (vallée de l'étrange), \emph{la scomparsa dell'autore}, \emph{il lavoro di bottega}, \emph{restio a}.

\end{indication}

\hypertarget{solution-8}{%
\subsubsection{Solution}\label{solution-8}}

Le conditionnel (\emph{potrebbe}, \emph{esisterebbe}) porte l'hypothèse initiale ; \emph{però}, \emph{Tuttavia}, \emph{Infine} ouvrent les trois temps suivants. La réticence à l'émotion est explicitement restituée (\emph{siamo restii a commuoverci}).

Di fronte alla massa di immagini prodotte dall'IA generativa, si potrebbe ridurla a uno strumento di illustrazione a basso costo. Non esisterebbe un'estetica artificiale: nel migliore dei casi, arte algoritmica; nel peggiore, una tecnica che trasferisce stili, plagia e produce abili collage, sorprendenti solo perché ricombinano forme note. Questa analisi ricorda però che gli artisti hanno sognato a lungo macchine creatrici e la scomparsa dell'autore, e invita a vedere nell'IA una produzione delegata, simile al lavoro di bottega. Tuttavia, sapere che queste opere provengono da una macchina ci fa cadere nella «valle perturbante»: non ne ricostruiamo l'intenzione e siamo restii a commuoverci. Infine, queste immagini non mostrano né il mondo né uno sguardo umano, ma rappresentazioni riorganizzate secondo regole proprie della macchina. Spetta all'artista e allo spettatore decidere che fare di questa verità estranea.

\textbf{(134 parole)}

\begin{methode}

Le pronom \emph{ne} et le pronominal remplacent une reprise lourde : \emph{non ne ricostruiamo l'intenzione} (l'intention de ces œuvres), \emph{commuoverci}. Ils font gagner des mots sans perdre de sens.

\end{methode}

\begin{piege}

\emph{Collage} est invariable en italien (\emph{abili collage}), comme la plupart des mots étrangers. Et \emph{restio} s'accorde : \emph{siamo restii}.

\end{piege}

\begin{jury}

Rapport 2026 (p.~10) : le jury rappelle « l'importance de reformuler les propos afin de se les approprier » et note qu'un élément a souvent été évacué : « l'affirmation de la réticence des humains à éprouver de l'émotion devant une production artistique produite par une IA générative ».

\end{jury}

\hypertarget{question-10-italien-ii.-ruxe9daction}{%
\subsection{Question 10 --- Italien, II. Rédaction}\label{question-10-italien-ii.-ruxe9daction}}

\begin{indication}{1}

\emph{In che senso} ne demande pas « dans quelle mesure » mais « en quel sens » : il faut dire \textbf{sur quoi} porte la révolution, et sur quoi elle ne porte pas.

\end{indication}

\begin{indication}{2}

Un exemple italien récent est un atout : l'encart rédigé par l'IA que \emph{Il Foglio} a mis en kiosque en mars 2025, la loi italienne sur l'IA adoptée en 2025.

\end{indication}

\begin{indication}{3}

Le jury apprécie les exemples originaux : chaque argument doit en avoir un.

\end{indication}

\hypertarget{solution-9}{%
\subsubsection{Solution}\label{solution-9}}

Plan : introduction sur \emph{Il Foglio AI} et reformulation de la question ; I. l'IA révolutionne les pratiques (musique, cinéma, jeu vidéo, démocratisation et ses coûts, la loi de 2025) ; II. une révolution technique plus qu'esthétique (photographie, expérience vécue, intention) ; conclusion.

Nel marzo 2025 il quotidiano Il Foglio ha messo in edicola un inserto interamente scritto dall'intelligenza artificiale. Questo esperimento dimostra che l'IA non riguarda più soltanto le immagini, ma tutte le forme di creazione. In che senso, però, si può parlare di rivoluzione?

Innanzitutto, l'IA rivoluziona le pratiche artistiche. Nella musica, alcuni programmi compongono ormai una canzone completa in pochi secondi; nel cinema e nei videogiochi, generano scenografie, voci e personaggi. Chiunque, anche senza alcuna formazione, può così diventare creatore. Questa democratizzazione ha però un costo: illustratori, doppiatori e musicisti temono per il loro lavoro, tanto che nel 2025 il Parlamento italiano ha approvato una legge che riconosce il diritto d'autore solo alle opere frutto dell'ingegno umano, anche se create con l'aiuto dell'IA.

Tuttavia, la rivoluzione è più tecnica che estetica. Già nell'Ottocento la fotografia sembrava condannare la pittura, ma in realtà ha liberato i pittori dall'obbligo di imitare la realtà. Allo stesso modo, l'IA potrebbe spingere gli artisti verso ciò che una macchina non sa fare: esprimere esperienze vissute. Senza intenzione umana, ogni immagine resta una semplice combinazione di dati.

In conclusione, secondo me l'IA sta rivoluzionando gli strumenti e il mercato dell'arte, ma non il suo senso, che dipende ancora dall'uomo.

\textbf{(203 parole)}

\begin{methode}

\emph{Già nell'Ottocento \ldots{} Allo stesso modo \ldots{}} : un précédent historique suivi d'une analogie donne une partie entière avec peu de mots.

\end{methode}

\begin{piege}

Les formes élidées (\emph{dall'intelligenza}, \emph{un'immagine}) comptent en général pour un mot. Si l'on craint une convention plus stricte, il faut vérifier que le texte reste dans la fourchette quand on les compte pour deux : ici 203 mots en les comptant pour un, 216 pour deux.

\end{piege}

\begin{jury}

Rapport 2026 (p.~10) : les arguments « ont toujours été pertinents », et les candidats ont souvent trouvé des exemples originaux. Le jury insiste sur l'importance de « structurer ses idées de manière à pouvoir les justifier à travers des exemples, même si le nombre de mots est limité ».

\end{jury}

\hypertarget{bilan}{%
\section{Bilan}\label{bilan}}

\textbf{À retenir pour toutes les langues.}

\begin{itemize}
\tightlist
\item
  Contraction : lire le texte français jusqu'à en dégager les quatre temps ; répartir les mots avant d'écrire ; supprimer les exemples et les noms propres ; reformuler sans rien ajouter ; marquer chaque bascule par un connecteur. Le dernier paragraphe, le plus abstrait, est celui que tous les jurys voient sacrifié.
\item
  Rédaction : répondre à la question posée (le degré ou le sens de la « révolution »), en deux parties annoncées par une phrase-thèse, avec des exemples précis de plusieurs arts, pris hors du texte et si possible dans l'aire de la langue.
\item
  Compter ses mots et les indiquer à la fin de chaque exercice. Viser le milieu de la fourchette laisse une marge si les conventions diffèrent (élisions, nombres, traits d'union).
\end{itemize}

\textbf{Grammaire à revoir, langue par langue.} Allemand : place du verbe, déclinaison, genre et pluriel, \emph{Konjunktiv I}. Anglais : modaux suivis de l'infinitif, indénombrables, \emph{AI} sans article et repris par \emph{it}, concession. Arabe : accord du verbe placé avant son sujet, duel, accord du nom après un nombre. Espagnol : conjugaisons, \emph{a} devant un complément humain, \emph{ser} / \emph{estar}, accentuation. Italien : accord des adjectifs, pronoms \emph{ne} et \emph{ci}, mots étrangers invariables.

\textbf{Situations qui se retrouvent dans d'autres sujets.} Contracter un texte théorique qui avance par concession (« certes \ldots{} mais »). Juger l'ampleur d'une rupture technique dans un domaine culturel. Mobiliser le précédent historique (la photographie face à la peinture) pour relativiser une nouveauté.

\end{document}
